\documentclass[a4paper,12pt]{article}

\usepackage[a4paper,margin=1in]{geometry}
\usepackage{enumitem}
\usepackage{listings}
\usepackage{xcolor}
\usepackage{hyperref}
\usepackage{graphicx}

\setlength{\parindent}{0pt}
\setlength{\parskip}{0.6em}

% Assembly style
\lstdefinestyle{AsmStyle}{
    language={[x86masm]Assembler},
    basicstyle=\ttfamily\footnotesize,
    keywordstyle=\color{blue},
    commentstyle=\color{green!40!black},
    numbers=left,
    numberstyle=\tiny,
    breaklines=true,
    frame=single
}

\title{Project 2 Report: \texttt{schedsim}\\
\large CMPE 230 Systems Programming -- Spring 2026}

\author{
Kemal Önal 
}

\date{\today}

\begin{document}

\maketitle

\section{Project Overview}

This program is a CPU scheduling simulator written in x86-64 GNU assembly. It reads a single line of input describing a scheduling algorithm and a set of processes, then simulates the scheduler clock cycle by clock cycle and prints the resulting execution timeline. The output is a string of characters where each character represents one clock cycle: a letter for the running process or \texttt{'X'} for an idle CPU.

The program supports five scheduling algorithms: First Come First Serve (FCFS), Shortest Job First (SJF), Shortest Remaining Time First (SRTF), Priority First (PF), and Round Robin (RR). Two of them are non-preemptive (FCFS, SJF) and three are preemptive (SRTF, PF, RR). Each algorithm has its own tie-breaking rules to ensure deterministic output.

The overall execution flow has three main phases:
\begin{enumerate}
    \item \textbf{Input:} Read the entire input line from stdin using the \texttt{sys\_read} syscall.
    \item \textbf{Parsing:} Walk through the input buffer character by character to identify the algorithm and parse all process descriptors into parallel arrays.
    \item \textbf{Simulation:} Run a clock-cycle loop that selects which process runs each cycle, writes the corresponding character to the output buffer, and updates the state. When all processes are done, write the output buffer to stdout and exit.
\end{enumerate}

\subsection{File Structure}

The code is split into eight assembly files to keep each part focused and readable:

\begin{itemize}
    \item \texttt{data.s} -- Defines all shared variables in the \texttt{.bss} section. Every other file references these symbols.
    \item \texttt{schedsim.s} -- Contains \texttt{\_start} (entry point), the main simulation loop, the algorithm dispatcher, \texttt{exec\_proc} (writes output), and the exit logic.
    \item \texttt{parser.s} -- Contains all input parsing: algorithm detection, token iteration, process descriptor parsing, and the \texttt{atoi} function.
    \item \texttt{fcfs.s} -- First Come First Serve scheduling logic.
    \item \texttt{sjf.s} -- Shortest Job First scheduling logic.
    \item \texttt{srtf.s} -- Shortest Remaining Time First scheduling logic.
    \item \texttt{pf.s} -- Priority First scheduling logic.
    \item \texttt{rr.s} -- Round Robin scheduling logic with queue operations.
\end{itemize}

All files are assembled individually with \texttt{as} and linked together with \texttt{ld}. Cross-file references (like \texttt{do\_fcfs} or \texttt{exec\_proc}) are resolved by the linker through \texttt{.global} directives. No C library is used.

\section{Parsing Logic}

\subsection{Input Processing Strategy}

The input is read into a 4096-byte buffer called \texttt{input\_buf} using \texttt{sys\_read} (syscall 0). After reading, we set \texttt{rsi} to the start of the buffer and begin parsing.

The parsing has two phases. In the first phase, we identify the algorithm by checking the first characters of the input. The mapping is:

\begin{itemize}
    \item \texttt{'F'} + \texttt{'C'} $\rightarrow$ FCFS (\texttt{algo\_type = 0}), skip 4 bytes
    \item \texttt{'S'} + \texttt{'J'} $\rightarrow$ SJF (\texttt{algo\_type = 1}), skip 3 bytes
    \item \texttt{'S'} + anything else $\rightarrow$ SRTF (\texttt{algo\_type = 2}), skip 4 bytes
    \item \texttt{'P'} $\rightarrow$ PF (\texttt{algo\_type = 3}), skip 2 bytes
    \item anything else $\rightarrow$ RR (\texttt{algo\_type = 4}), skip 2 bytes
\end{itemize}

This works because the specification guarantees valid input, so we only need to distinguish between the five algorithm names without checking every character.

In the second phase, we iterate over the remaining tokens. Tokens are separated by single spaces (guaranteed by spec). Each token is either a process descriptor like \texttt{A-3-1} or, in the case of RR, the quantum value at the end. We distinguish them by checking the first character: if the algorithm is RR and the character is a digit, it is the quantum; otherwise, it is a process descriptor.

\subsection{Process Descriptor Parsing}

A process descriptor has the format \texttt{ID-BurstTime[-ArrivalTime[-Priority]]} depending on the algorithm. The parsing steps for one descriptor are:

\begin{enumerate}
    \item Read the ID character at \texttt{(\%rsi)} and store it in \texttt{proc\_id[count]}.
    \item Store the current count as the input order in \texttt{proc\_order[count]}.
    \item Skip past the ID and the hyphen with two \texttt{inc \%rsi} instructions.
    \item Call \texttt{atoi} to parse the burst time. Store the result in both \texttt{proc\_burst[]} and \texttt{proc\_rem[]}.
    \item If the algorithm is FCFS, SRTF, or PF: skip the hyphen and call \texttt{atoi} again for arrival time.
    \item If the algorithm is PF: skip another hyphen and call \texttt{atoi} for priority.
\end{enumerate}

We do not explicitly check for the hyphen character. Instead, we use \texttt{inc \%rsi} to skip over it. This works because \texttt{atoi} stops at the first non-digit character (the hyphen), so after \texttt{atoi} returns, \texttt{rsi} points exactly at the hyphen. One \texttt{inc} moves past it to the next digit.

Here is the relevant code for parsing a process descriptor:

\begin{lstlisting}[style=AsmStyle]
parse_proc:
    movl proc_count(%rip), %ecx    ; ecx = current index
    movzb (%rsi), %eax             ; read ID character
    lea proc_id(%rip), %rdi
    movb %al, (%rdi, %rcx, 1)      ; proc_id[i] = ID
    lea proc_order(%rip), %rdi
    movb %cl, (%rdi, %rcx, 1)      ; proc_order[i] = i
    inc %rsi                        ; skip ID
    inc %rsi                        ; skip '-'
    push %rcx
    call atoi                       ; parse burst time
    pop %rcx
    lea proc_burst(%rip), %rdi
    movl %eax, (%rdi, %rcx, 4)     ; burst[i] = value
    lea proc_rem(%rip), %rdi
    movl %eax, (%rdi, %rcx, 4)     ; rem[i] = value
\end{lstlisting}

\subsection{The atoi Function}

Since we cannot use the C library, we wrote our own \texttt{atoi} function to convert ASCII digit strings to integers. It takes the current position in \texttt{rsi} and returns the integer value in \texttt{eax}. The algorithm is simple: start with \texttt{eax = 0}, and for each digit character, compute \texttt{eax = eax * 10 + (digit - '0')}.

\begin{lstlisting}[style=AsmStyle]
atoi:
    xor %eax, %eax             ; result = 0
atoi_loop:
    movzb (%rsi), %ecx         ; load current byte
    cmp $'0', %ecx
    jl atoi_end                ; not a digit -> done
    cmp $'9', %ecx
    jg atoi_end                ; not a digit -> done
    sub $'0', %ecx             ; convert ASCII to value
    imul $10, %eax             ; result *= 10
    add %ecx, %eax             ; result += digit
    inc %rsi                   ; advance pointer
    jmp atoi_loop
atoi_end:
    ret
\end{lstlisting}

The function uses \texttt{call}/\texttt{ret} so it can be invoked from multiple places. One important detail is that \texttt{atoi} overwrites \texttt{ecx}. During parsing, \texttt{ecx} holds the current process index, so we save it with \texttt{push \%rcx} before calling \texttt{atoi} and restore it with \texttt{pop \%rcx} afterward.

For the RR quantum, we call \texttt{atoi} without saving \texttt{ecx} because after parsing the quantum we jump directly to \texttt{parse\_done} and \texttt{ecx} is no longer needed.

\subsection{Parsing Example}

To illustrate how parsing works, consider the input \texttt{SRTF A-4-0 B-2-1}.

\begin{enumerate}
    \item \texttt{rsi} points to \texttt{'S'}. First char is \texttt{'S'}, second is \texttt{'R'} (not \texttt{'J'}) $\rightarrow$ SRTF. \texttt{algo\_type = 2}. Skip 4 bytes.
    \item \texttt{rsi} points to \texttt{' '}. Skip space. Now points to \texttt{'A'}.
    \item Parse process 0: ID = \texttt{'A'}, skip \texttt{'A'} and \texttt{'-'}, \texttt{atoi} reads \texttt{"4"} $\rightarrow$ burst=4. Skip \texttt{'-'}, \texttt{atoi} reads \texttt{"0"} $\rightarrow$ arrival=0.
    \item \texttt{rsi} points to \texttt{' '}. Skip space. Now points to \texttt{'B'}.
    \item Parse process 1: ID = \texttt{'B'}, burst=2, arrival=1.
    \item \texttt{rsi} hits newline $\rightarrow$ jump to \texttt{parse\_done}.
\end{enumerate}

\section{State Representation (Memory Layout)}

\subsection{Process Representation}

We chose to represent processes using separate parallel arrays rather than a struct-of-arrays layout. Each attribute has its own array in the \texttt{.bss} section. This makes it easy to access any single attribute across all processes using a loop with scaled indexing.

The arrays and their sizes:

\begin{center}
\begin{tabular}{|l|l|l|l|}
\hline
\textbf{Array} & \textbf{Element Size} & \textbf{Total Size} & \textbf{Purpose} \\
\hline
\texttt{proc\_id[]} & 1 byte & 16 bytes & Uppercase letter ID \\
\texttt{proc\_order[]} & 1 byte & 16 bytes & Input index for tie-breaking \\
\texttt{proc\_burst[]} & 4 bytes & 64 bytes & Original burst time \\
\texttt{proc\_arrival[]} & 4 bytes & 64 bytes & Arrival time (default 0) \\
\texttt{proc\_rem[]} & 4 bytes & 64 bytes & Remaining burst time \\
\texttt{proc\_prio[]} & 4 bytes & 64 bytes & Priority (PF only, default 0) \\
\hline
\end{tabular}
\end{center}

The sizes are allocated for up to 16 processes, which is more than the specification maximum of 10. We access elements using x86-64 scaled indexing. For byte arrays we use \texttt{(\%rdi, \%rcx, 1)}, and for 4-byte arrays we use \texttt{(\%rdi, \%rcx, 4)}. The base address is loaded into \texttt{rdi} with \texttt{lea}, and the process index is in \texttt{rcx} or \texttt{rax} depending on the context.

For example, to read the remaining time of process \texttt{i}:

\begin{lstlisting}[style=AsmStyle]
    lea proc_rem(%rip), %rdi    ; rdi = base address of proc_rem
    movl (%rdi, %rax, 4), %edx  ; edx = proc_rem[eax]
\end{lstlisting}

The \texttt{proc\_burst} and \texttt{proc\_rem} arrays are initialized to the same value during parsing. \texttt{proc\_burst} stays constant throughout the simulation (used by SJF for selection), while \texttt{proc\_rem} is decremented each time a process runs.

\subsection{Global State}

The simulation state is tracked by several 4-byte global variables:

\begin{center}
\begin{tabular}{|l|l|l|}
\hline
\textbf{Variable} & \textbf{Initial Value} & \textbf{Purpose} \\
\hline
\texttt{current\_time} & 0 & Current clock cycle \\
\texttt{out\_len} & 0 & Number of characters written to output \\
\texttt{running\_proc} & -1 & Index of running process (-1 = idle) \\
\texttt{proc\_count} & (parsed) & Total number of processes \\
\texttt{algo\_type} & (parsed) & Algorithm code (0--4) \\
\hline
\end{tabular}
\end{center}

The output is built in a 4096-byte buffer called \texttt{output\_buf}. Each cycle, one character is written at position \texttt{out\_len} and the counter is incremented. At the end, a newline is appended and the entire buffer is written to stdout with a single \texttt{sys\_write} call.

\subsection{Round Robin Queue}

Round Robin uses a linear FIFO queue stored in \texttt{rr\_queue[]}, a 4096-byte array. Each element is 1 byte representing a process index. The queue is managed with two indices:

\begin{itemize}
    \item \texttt{rr\_head} -- index of the next process to dequeue.
    \item \texttt{rr\_tail} -- index of the next empty slot for enqueue.
\end{itemize}

When \texttt{rr\_head == rr\_tail}, the queue is empty. Both indices only increase and never wrap around. This is a linear queue rather than a circular one, but it works because the total number of enqueue operations in any valid input is bounded by the output length (at most 1024 characters), which is well under the 4096-byte buffer.

There is also \texttt{rr\_quantum\_left}, which counts how many cycles remain in the current quantum slot. It is set to the quantum value when a process is dequeued and decremented each cycle.

\subsection{Mandatory Diagram (Hand-Drawn)}

\includegraphics[width=\textwidth]{diagram.jpg}

\section{Register Usage Strategy}

The register usage follows a consistent pattern across all scheduling algorithm files:

\begin{itemize}
    \item \texttt{rsi} -- input buffer pointer during parsing, also used to load \texttt{proc\_order} base address during scheduling.
    \item \texttt{rdi} -- used to load base addresses of data arrays (e.g., \texttt{lea proc\_rem(\%rip), \%rdi}).
    \item \texttt{eax} -- loop counter during process iteration, also used for syscall numbers.
    \item \texttt{ecx} -- process count (loop bound), also process index during parsing.
    \item \texttt{edx} -- holds the current candidate's comparison value (arrival, burst, remaining, or priority depending on the algorithm).
    \item \texttt{r9d} -- index of the best candidate found so far.
    \item \texttt{r10d} -- best value of the primary selection criterion found so far.
    \item \texttt{r11d} -- best input order found so far (for tie-breaking).
    \item \texttt{r12d} -- best remaining time (used only in PF for the second-level tie-break).
    \item \texttt{r13d} -- current process remaining time (used only in PF).
    \item \texttt{cl} -- holds the output character (process ID letter or \texttt{'X'} for idle).
\end{itemize}

We do not follow a strict caller-saved / callee-saved convention since we do not use external functions. The only internal function call is \texttt{atoi}, and we manually save \texttt{rcx} with \texttt{push}/\texttt{pop} around it because \texttt{atoi} uses \texttt{ecx} internally.

\section{Scheduling Logic}

\subsection{Clock Cycle Simulation}

The simulation runs in a main loop (\texttt{sim\_loop}). At the start of each iteration, we check if all processes have a remaining time of 0. If so, and there is no active RR quantum to pad, we exit. Otherwise, we dispatch to the appropriate algorithm handler based on \texttt{algo\_type}.

Each algorithm handler selects a process (or decides the CPU is idle) and stores the result in \texttt{running\_proc}. Then control jumps to \texttt{exec\_proc}, which loads the process ID from \texttt{proc\_id[]}, writes it to \texttt{output\_buf[out\_len]}, decrements the remaining time, increments \texttt{out\_len} and \texttt{current\_time}, and loops back to \texttt{sim\_loop}. If \texttt{running\_proc} is -1, it writes \texttt{'X'} instead.

\subsection{FCFS -- First Come First Serve}

FCFS is non-preemptive. If \texttt{running\_proc} is not -1 and the running process still has remaining time, it continues without interruption. Only when the current process finishes (remaining reaches 0) do we search for a new one.

The selection scans all processes and picks the one with the smallest arrival time that has already arrived (\texttt{arrival <= current\_time}) and still has remaining time. If two processes have the same arrival time, we pick the one that appeared first in the input by comparing their \texttt{proc\_order} values.

For example, with input \texttt{FCFS T-3-1 B-2-2 A-3-1}: at time 1, both T (arrival=1, order=0) and A (arrival=1, order=2) arrive. T has a lower order, so T runs first.

\subsection{SJF -- Shortest Job First}

SJF has the same structure as FCFS but with two differences: there is no arrival time check (all processes arrive at time 0), and the selection criterion is the original burst time (\texttt{proc\_burst}) instead of arrival time. This means we always compare the original burst, not the remaining time, even after a process has started.

Since SJF is non-preemptive, a running process continues until it finishes. Ties are broken by input order.

\subsection{SRTF -- Shortest Remaining Time First}

SRTF is preemptive. Unlike FCFS and SJF, there is no ``continue running'' check. Every single cycle, we re-scan all processes and pick the arrived one with the smallest remaining time. This means if a new process arrives with a shorter remaining time than the currently running one, it automatically takes over on the next cycle.

The selection uses \texttt{proc\_rem} (remaining time) instead of \texttt{proc\_burst}, and it checks \texttt{proc\_arrival <= current\_time} to only consider arrived processes. Ties are broken by input order.

\subsection{PF -- Priority First}

PF is preemptive and uses a three-level comparison. Every cycle, we scan all arrived processes and select by:

\begin{enumerate}
    \item Lowest priority number (lower = higher priority).
    \item If priorities are equal: shortest remaining time (SRTF behavior).
    \item If both are equal: earliest input order.
\end{enumerate}

This requires two extra registers compared to the other algorithms: \texttt{r12d} tracks the best remaining time found so far, and \texttt{r13d} holds the current candidate's remaining time. The comparison chain is: check priority with \texttt{cmp \%r10d, \%edx}, then remaining with \texttt{cmp \%r12d, \%r13d}, then order with \texttt{cmp \%r11d, \%r8d}.

\subsection{RR -- Round Robin}

RR is the most structurally different algorithm. It does not use \texttt{exec\_proc} at all and handles its own output writing. The flow is:

\begin{enumerate}
    \item At the start, all processes are enqueued in input order.
    \item If no process is running (\texttt{running\_proc == -1}), dequeue the front process and set \texttt{quantum\_left = Q}.
    \item Each cycle: decrement \texttt{quantum\_left}. If \texttt{remaining > 0}, output the process ID and decrement remaining. If \texttt{remaining == 0}, output \texttt{'X'} as padding.
    \item When \texttt{quantum\_left} reaches 0: if the process still has remaining time, re-enqueue it at the back. Set \texttt{running\_proc = -1} and go back to step 2.
\end{enumerate}

The padding behavior is important: if a process finishes before the quantum expires, the remaining quantum cycles are filled with \texttt{'X'}. The CPU does not switch to the next process mid-quantum.

One subtle detail is the quantum counting. When a process is dequeued, \texttt{quantum\_left} is set to Q, then \texttt{jmp do\_rr} brings us back to the top where \texttt{decl quantum\_left} runs immediately. So the first cycle decrements Q to Q-1, the second to Q-2, and so on. The check \texttt{cmpl \$0, quantum\_left; jg continue} means the process runs for exactly Q cycles before the quantum expires. We verified this by tracing with small quantum values on paper.

\section{Control Flow Design}

The program is structured as a sequence of phases connected by jumps:

\begin{enumerate}
    \item \texttt{\_start} reads input and jumps to \texttt{skip\_spc} in the parser.
    \item The parser identifies the algorithm, parses all process tokens, and jumps to \texttt{parse\_done} in the main file.
    \item \texttt{parse\_done} initializes the simulation state and falls through to \texttt{sim\_loop}.
    \item \texttt{sim\_loop} checks the termination condition and jumps to \texttt{sim\_continue}, which dispatches to one of the five algorithm handlers (\texttt{do\_fcfs}, \texttt{do\_sjf}, etc.).
    \item Each handler jumps to \texttt{exec\_proc} (or handles output internally for RR), which writes output and jumps back to \texttt{sim\_loop}.
    \item When all processes finish, \texttt{sim\_exit} writes the output buffer to stdout and terminates with \texttt{sys\_exit}.
\end{enumerate}

All scheduling algorithms use the same loop pattern: initialize sentinel values (\texttt{r9d = -1}, \texttt{r10d = 0x7FFFFFFF}), iterate over all processes with a counter in \texttt{eax}, compare against the current best, and update if better. Conditional jumps (\texttt{jl}, \texttt{jg}, \texttt{je}, \texttt{jle}, \texttt{jge}) are used for all comparisons, which are signed comparisons.

\section{Representative Code Snippets}

\subsection{SRTF Selection Loop}

\begin{lstlisting}[style=AsmStyle]
srtf_loop:
    cmp %ecx, %eax
    jge srtf_end
    lea proc_rem(%rip), %rdi
    movl (%rdi, %rax, 4), %edx     ; edx = remaining[i]
    test %edx, %edx
    jle srtf_next                   ; skip if done
    lea proc_arrival(%rip), %rdi
    movl (%rdi, %rax, 4), %esi
    cmpl current_time(%rip), %esi
    jg srtf_next                    ; skip if not arrived

    cmp %r10d, %edx
    jl srtf_set                     ; shorter remaining -> new best
    jg srtf_next
    lea proc_order(%rip), %rsi
    movzb (%rsi, %rax, 1), %r8d
    cmp %r11d, %r8d
    jge srtf_next                   ; same remaining, later input -> skip
srtf_set:
    movl %eax, %r9d                 ; best index = i
    movl %edx, %r10d                ; best remaining = remaining[i]
    lea proc_order(%rip), %rsi
    movzb (%rsi, %rax, 1), %r11d    ; best order = order[i]
\end{lstlisting}

This is the core selection loop for SRTF. It iterates over all processes (counter in \texttt{eax}, bound in \texttt{ecx}). For each process, it first loads the remaining time into \texttt{edx} and checks if it is positive. Then it loads the arrival time into \texttt{esi} and checks if the process has arrived (\texttt{arrival <= current\_time}). If both conditions pass, it compares the remaining time against the current best in \texttt{r10d}. If equal, it falls through to the tie-breaking check using input order in \texttt{r11d}. The best candidate index is tracked in \texttt{r9d}. After the loop, \texttt{r9d} is stored into \texttt{running\_proc}.

All four non-RR algorithms (FCFS, SJF, SRTF, PF) follow this same pattern. The only differences are which arrays they compare and how many levels of tie-breaking they use.

\subsection{exec\_proc -- Writing Output}

\begin{lstlisting}[style=AsmStyle]
exec_proc:
    movl running_proc(%rip), %eax
    cmpl $-1, %eax
    jne exec_valid
    movb $'X', %cl               ; no process -> idle
    jmp exec_write
exec_valid:
    lea proc_id(%rip), %rdi
    movb (%rdi, %rax, 1), %cl    ; cl = process ID letter
    lea proc_rem(%rip), %rdi
    decl (%rdi, %rax, 4)         ; remaining[i]--
exec_write:
    lea output_buf(%rip), %rdi
    movl out_len(%rip), %edx
    movb %cl, (%rdi, %rdx, 1)    ; output_buf[out_len] = cl
    incl out_len(%rip)
    incl current_time(%rip)
    jmp sim_loop
\end{lstlisting}

This function is shared by FCFS, SJF, SRTF, and PF. It checks if \texttt{running\_proc} is -1, and if so, writes \texttt{'X'} for an idle cycle. Otherwise, it loads the process ID character into \texttt{cl}, decrements the remaining time, writes the character to \texttt{output\_buf} at position \texttt{out\_len}, then increments both \texttt{out\_len} and \texttt{current\_time} before looping back.

\section{Testing and Debugging}

\subsection{Testing Strategy}

We tested the program using the provided test cases through \texttt{make testcases} and \texttt{make grade}. The Makefile runs each test case by piping the input file into the executable and capturing the output. The grader script then compares our output against the expected output using exact string matching.

We also manually traced several edge case inputs on paper before running them to verify our expectations matched the actual output. This was especially helpful for RR, where the interaction between quantum, padding, and re-enqueue can be confusing.

\subsection{Debugging Approach}

Since we cannot use \texttt{gdb} easily with a statically linked, no-libc assembly program, we relied on two main debugging techniques:

\begin{enumerate}
    \item \textbf{Print debugging:} We temporarily added \texttt{sys\_write} calls at specific points in the code to print register values or marker characters to stdout. For example, we would write the value of \texttt{running\_proc} after each selection to see if the right process was being chosen. These debug prints were removed from the final submission.
    \item \textbf{Paper tracing:} We traced the execution manually on paper for small inputs, writing down register values at each step of the simulation loop. This was particularly useful for finding off-by-one errors in the RR quantum counting logic and for verifying tie-breaking behavior.
\end{enumerate}

One specific bug we found through paper tracing was in the RR quantum counting. Initially, we were not sure if \texttt{decl quantum\_left} happening immediately after dequeue would cause the process to run for Q-1 cycles instead of Q. We traced through with Q=2 and confirmed that the \texttt{jg} check (strictly greater than 0) makes it work correctly for exactly Q cycles.

\subsection{Edge Cases}

The edge cases we focused on:

\begin{itemize}
    \item \textbf{Tie-breaking:} Multiple processes arriving at the same time with the same burst or priority. For example, \texttt{FCFS T-3-1 B-2-2 A-3-1} has T and A arriving at the same time. We verified that T runs first because it appears earlier in the input.
    \item \textbf{Idle CPU:} When no process has arrived yet, the output should contain \texttt{'X'}. For example, \texttt{FCFS A-2-3} should produce \texttt{XXXAA} because A does not arrive until cycle 3.
    \item \textbf{RR padding:} When a process finishes before its quantum expires, the remaining quantum cycles should be padded with \texttt{'X'}. For example, with \texttt{RR A-1 3}, process A finishes in 1 cycle but holds the quantum for 3, producing \texttt{AXX}.
    \item \textbf{Preemption:} In SRTF, a newly arrived process with a shorter remaining time should immediately take over. For example, \texttt{SRTF A-4-0 B-2-1} produces \texttt{ABBAAAA} because B (rem=2) preempts A (rem=3) at cycle 1.
    \item \textbf{Single process:} We tested each algorithm with just one process to make sure the base case works.
    \item \textbf{All same burst:} In SJF, if all processes have the same burst time, they should run in input order.
\end{itemize}

\subsection{Sample Test Trace}

To show how we verified correctness, here is a step-by-step trace for \texttt{RR A-3 B-4 C-2 2}:

\begin{center}
\begin{tabular}{|c|c|c|c|c|c|}
\hline
\textbf{Turn} & \textbf{Proc} & \textbf{Rem Before} & \textbf{Cycles} & \textbf{Output} & \textbf{Rem After} \\
\hline
1 & A & 3 & 2 & AA & 1 \\
2 & B & 4 & 2 & BB & 2 \\
3 & C & 2 & 2 & CC & 0 \\
4 & A & 1 & 1+1 pad & AX & 0 \\
5 & B & 2 & 2 & BB & 0 \\
\hline
\end{tabular}
\end{center}

Final output: \texttt{AABBCCAXBB}. This matches the expected output from the specification.

\section{Course Concepts}

\subsection{Memory and Data Layout}

The program uses the \texttt{.bss} section for all data, which is zero-initialized by the OS. We use parallel arrays with scaled indexing to access process attributes. This maps directly to how arrays work at the hardware level: a base address plus an offset computed from the index and element size. The \texttt{lea} instruction loads the base address using RIP-relative addressing, and the scaled index form \texttt{(\%rdi, \%rax, 4)} computes the final address in a single instruction.

One design decision was storing process IDs and input order as 1-byte values and burst/arrival/remaining/priority as 4-byte values. This means we use scale factor 1 for IDs and scale factor 4 for numeric attributes. Since x86-64 supports both natively, this does not add complexity.

\subsection{Registers and Low-Level Computation}

All scheduling decisions are made using register-level comparisons and conditional jumps. In a high-level language, selecting the best process would be a simple loop with an if statement. In assembly, this becomes a series of \texttt{cmp} instructions followed by conditional jumps (\texttt{jl}, \texttt{jg}, \texttt{je}). The flags register implicitly carries the comparison result between the \texttt{cmp} and the jump.

We use \texttt{jl}/\texttt{jg} (signed comparisons) rather than \texttt{jb}/\texttt{ja} (unsigned). This is correct because all values are non-negative integers that fit in 32 bits. The sentinel value \texttt{0x7FFFFFFF} is the largest positive signed 32-bit integer, so any valid value will compare as less than it.

\subsection{Control Flow and Branching}

Loops are implemented with labels and backward jumps. The typical pattern is:

\begin{lstlisting}[style=AsmStyle]
    movl $0, %eax           ; i = 0
loop:
    cmp %ecx, %eax          ; i >= count?
    jge done                ; if so, exit loop
    ; ... loop body ...
    inc %eax                ; i++
    jmp loop                ; repeat
done:
\end{lstlisting}

This is the assembly equivalent of a \texttt{for (int i = 0; i < count; i++)} loop. All five scheduling algorithms and the termination check use this pattern.

Conditional logic (if/else) is implemented using \texttt{cmp} followed by a conditional jump that skips over the "then" body. For multi-level comparisons like in PF (priority, then remaining, then order), we chain multiple \texttt{cmp}/\texttt{jl}/\texttt{jg} pairs, where \texttt{jl} jumps to the ``set new best'' label and \texttt{jg} jumps to the ``skip'' label. If neither fires, we fall through to the next comparison level.

\subsection{Parsing Without Libraries}

Since we cannot use \texttt{libc}, all string processing is done manually. The \texttt{atoi} function converts ASCII digit strings to integers using the classic multiply-and-add approach: \texttt{result = result * 10 + (char - '0')}. Token boundaries are handled by advancing the pointer past delimiters with \texttt{inc \%rsi}.

This is fundamentally different from high-level parsing where you might use \texttt{strtok} or \texttt{sscanf}. In assembly, the pointer is just a register that we increment manually, and we check each byte with \texttt{movzb} and \texttt{cmp}. There is no concept of ``returning a string'' -- the pointer just moves forward as we consume characters.

\subsection{Tradeoffs}

We chose simplicity over efficiency in several places:

\begin{itemize}
    \item \textbf{Linear scan vs. sorted structure:} Every scheduling algorithm scans all processes each cycle (O(n) per cycle) instead of maintaining a sorted list or priority queue. With a maximum of 10 processes, this is fast enough.
    \item \textbf{Linear queue vs. circular queue:} The RR queue only grows and never wraps around. This wastes memory but avoids the complexity of modular arithmetic. With at most 1024 output characters, the queue never exceeds the 4096-byte buffer.
    \item \textbf{No delimiter validation:} We skip hyphens with blind \texttt{inc \%rsi} instead of checking that the character is actually a hyphen. This saves instructions but would break on malformed input. The specification guarantees valid input, so this is acceptable.
    \item \textbf{Signed comparisons:} We use signed \texttt{jl}/\texttt{jg} everywhere. This works because all values are non-negative and fit in 31 bits, but it would break if values exceeded \texttt{0x7FFFFFFF}.
\end{itemize}

\section{Conclusion}

The strongest design decision was decomposing the code into separate files per algorithm. This made it easy to work on each algorithm independently without worrying about breaking the others. It also makes the code easier to read, since each file is focused on one thing.

Another good decision was using the same register convention across all algorithms (\texttt{r9d} for best index, \texttt{r10d} for best value, \texttt{r11d} for best order). This means once you understand how one algorithm works, the others follow the same pattern.

The biggest challenge was getting the RR quantum counting right. The dequeue logic sets \texttt{quantum\_left} to Q, then immediately jumps back to \texttt{do\_rr} where it gets decremented. We had to carefully trace through the code on paper to verify that the process runs for exactly Q cycles and not Q-1. The key insight is that \texttt{jg} (jump if greater than 0) means the quantum expires when it reaches 0, not when it reaches -1.

Another challenge was making sure \texttt{atoi} does not corrupt the process index register. We solved this with \texttt{push}/\texttt{pop} around each call, but it took some debugging to realize that \texttt{atoi} was the one clobbering \texttt{ecx}.

Possible improvements:
\begin{itemize}
    \item Making the RR queue circular for robustness with longer outputs.
    \item Adding explicit hyphen checks during parsing.
    \item Using a proper calling convention for \texttt{atoi} (callee-saved registers) instead of manual \texttt{push}/\texttt{pop}.
\end{itemize}

\section*{AI Usage Statement}

We used ChatGPT occasionally during the project. It was mainly used as an instructor to clarify x86-64 conventions like syscall numbers and register usage rules. In a few cases, we had it help with implementing simple utility functions such as the \texttt{atoi} routine. We also had it review the final code to check for potential bugs or edge case issues. The overall design, scheduling logic, and program structure were done by us. We understand all parts of the code and can explain them.

\end{document}
